% !TEX spellcheck = en_US



% ~~~~~~~~~~~~~~~~~~~~~~~~~~~~~~~~~~~~~~~ V
% (NC) 2026 Haluk Bingol
% github.com/halukbingol/LaTeX-Templates
%
% Licensees may copy, distribute, display, and perform the work and make derivative works 
% and remixes based on it only for non-commercial purposes. 
% ~~~~~~~~~~~~~~~~~~~~~~~~~~~~~~~~~~~~~~~ A




% =====================================================================
%  Stacks -- Learning by Example
%  ---------------------------------------------------------------------
%  Built on hblib/Template/hblibTemplate.tex with the libraries in ../../hblib/.
%  Programs and their real outputs are in codeoutput/:
%     codeoutput/<Name>.java  a program (classes without main() are only compiled)
%     codeoutput/<Name>.sh    a terminal session, run in a copy of codeoutput/<Name>/
%     codeoutput/<Name>.txt   what it printed, made by:  bash run_examples.sh
% =====================================================================

\documentclass[aspectratio=169,10pt,compress]{beamer}

% ~~~~~~~~~~~~~~~~~~~~~~~~~~~~~~~~~~~~~~~~~~~~~~~~~~~~~~~~~~~~~~~~~ V hblib
	\usepackage{../../hblib/hblibPresentation} % lib presentation: theme, i/N, \hSection, algorithm2e, ...
	\usepackage{../../hblib/hblibCodeoutput} % lib code + output: \lstcode, \lstterminal, ...
	\usepackage{../../hblib/hblibDatastructures} % lib data structures: \hString, \hStack, \hQueue
% ~~~~~~~~~~~~~~~~~~~~~~~~~~~~~~~~~~~~~~~~~~~~~~~~~~~~~~~~~~~~~~~~~ A hblib


% xcolor, array (and the L{width} column), listings, tikz and algorithm2e come from the hblib libraries
\usepackage[T1]{fontenc}
\usepackage{lmodern}
\usepackage{booktabs}

\newcommand{\win}{\textbf{Windows}}
\newcommand{\mac}{\textbf{macOS}}

% algorithm2e set-up (the package is loaded by hblibPresentation)
\SetAlCapFnt{\small}
\SetAlCapNameFnt{\small}
\SetKwProg{Fn}{function}{}{end}
\SetKw{KwError}{error}

% ---------------------------------------------------------------------
\title{Stacks}
\subtitle{Learning by Example}
\author[Bingol]{Haluk O. Bingol}

\institute[FCIS]{
    Faculty of Computer and Information Sciences\\
    Yeditepe University
}
\date{
	{\tiny \hbTimeStamp}
}

\begin{document}




% ~~~~~~~~~~~~~~~~~~~~~~~~~~~~~~~~~~~~~~ V frm
\begin{frame}[t,fragile]
  \titlepage
\end{frame}
% ~~~~~~~~~~~~~~~~~~~~~~~~~~~~~~~~~~~~~~ A frm




% ~~~~~~~~~~~~~~~~~~~~~~~~~~~~~~~~~~~~~~ V frm
\begin{frame}[t,fragile] \frametitle{Outline}
  \begin{columns}[T]
    \begin{column}{0.48\textwidth}
      \tableofcontents[sections={1-3}, hideallsubsections]
    \end{column}
    \begin{column}{0.48\textwidth}
      \tableofcontents[sections={4-6}, hideallsubsections]
    \end{column}
  \end{columns}
\end{frame}
% ~~~~~~~~~~~~~~~~~~~~~~~~~~~~~~~~~~~~~~ A frm

% ~~~~~~~~~~~~~~~~~~~~~~~~~~~~~~~~~~~~~~ V
\hSection[Setup]{How to Use These Examples}
% ~~~~~~~~~~~~~~~~~~~~~~~~~~~~~~~~~~~~~~ A




% ~~~~~~~~~~~~~~~~~~~~~~~~~~~~~~~~~~~~~~ V frm
\begin{frame}[t,fragile] \frametitle{Running the examples}
  \small Every program on these slides is a file in \texttt{codeoutput/}; the output shown is what it really printed (Java 21).
  \begin{columns}[T]
    \begin{column}{0.48\textwidth}
      \begin{block}{Windows (Command Prompt)}


% +++++++++++++++++++++++++++++++++++++++ V lst
%: -Lst.
\begin{lstlisting}[style=codesnippet, language={}, basicstyle=\ttfamily\scriptsize]
cd stack-2-examples\codeoutput
javac PushPop.java
java PushPop
\end{lstlisting}
% +++++++++++++++++++++++++++++++++++++++ A lst


      \end{block}
    \end{column}
    \begin{column}{0.48\textwidth}
      \begin{block}{macOS (Terminal)}


% +++++++++++++++++++++++++++++++++++++++ V lst
%: -Lst.
\begin{lstlisting}[style=codesnippet, language=bash, basicstyle=\ttfamily\scriptsize]
cd stack-2-examples/codeoutput
javac PushPop.java
java PushPop
\end{lstlisting}
% +++++++++++++++++++++++++++++++++++++++ A lst


      \end{block}
    \end{column}
  \end{columns}

  \begin{itemize}\small
    \item \texttt{javac PushPop.java} also compiles \texttt{IntStack.java}: it finds the class it needs in the same folder.
    \item \texttt{bash run\_examples.sh} reruns every program (\mac{}: Terminal; \win: Git Bash).
    \item Code first, output on the next click: \hRemark{predict the output before you click.}
  \end{itemize}
\end{frame}
% ~~~~~~~~~~~~~~~~~~~~~~~~~~~~~~~~~~~~~~ A frm




% ~~~~~~~~~~~~~~~~~~~~~~~~~~~~~~~~~~~~~~ V frm
\begin{frame}[t,fragile] \frametitle{What it looks like in a terminal}
  \lstterminal[basicstyle=\ttfamily\footnotesize]{codeoutput/Run.txt}

  \vspace{0.1em}
  \small
  Two \texttt{.class} files: \texttt{javac} compiled \texttt{IntStack} because \texttt{PushPop} uses it.
\end{frame}
% ~~~~~~~~~~~~~~~~~~~~~~~~~~~~~~~~~~~~~~ A frm

% ~~~~~~~~~~~~~~~~~~~~~~~~~~~~~~~~~~~~~~ V
\hSection[IntStack]{A Stack of int in an Array}
% ~~~~~~~~~~~~~~~~~~~~~~~~~~~~~~~~~~~~~~ A

% ~~~~~~~~~~~~~~~~~~~~~~~~~~~~~~~~~~~~~~ V
\subsection[Class]{The class}
% ~~~~~~~~~~~~~~~~~~~~~~~~~~~~~~~~~~~~~~ A




% ~~~~~~~~~~~~~~~~~~~~~~~~~~~~~~~~~~~~~~ V frm
\begin{frame}[t,fragile] \frametitle{\texttt{IntStack}: a fixed array and \texttt{top}}
  \begin{columns}[T]
    \begin{column}{0.49\textwidth}
      \lstcode[basicstyle=\ttfamily\tiny, lastline=19]{codeoutput/IntStack.java}{Java}
    \end{column}
    \begin{column}{0.49\textwidth}
      \lstcode[basicstyle=\ttfamily\tiny, firstline=20, firstnumber=20]{codeoutput/IntStack.java}{Java}
      \begin{center}\begin{tikzpicture}\hStack{0}{5}{2}{10,20,30}[0.45][0.45]\end{tikzpicture}\\ \small \texttt{top = 2}, \texttt{size() = 3}\end{center}
    \end{column}
  \end{columns}
  \small \texttt{a[++top] = x}: move \texttt{top} up, then store. \texttt{a[top--]}: read, then move \texttt{top} down.
  Every method checks before it touches the array; \texttt{toString} prints from the bottom to the top.
\end{frame}
% ~~~~~~~~~~~~~~~~~~~~~~~~~~~~~~~~~~~~~~ A frm

% ~~~~~~~~~~~~~~~~~~~~~~~~~~~~~~~~~~~~~~ V
\subsection[Use]{Using it}
% ~~~~~~~~~~~~~~~~~~~~~~~~~~~~~~~~~~~~~~ A




% ~~~~~~~~~~~~~~~~~~~~~~~~~~~~~~~~~~~~~~ V frm
\begin{frame}[t,fragile] \frametitle{Push, peek, pop}
  \lstcoderightstep[0.58][basicstyle=\ttfamily\tiny][basicstyle=\ttfamily\scriptsize]{codeoutput/PushPop.java}{Java}{codeoutput/PushPop.txt}

  \vspace{0.2em}
  \small \texttt{pop} returns the elements in the \textbf{reverse} order of \texttt{push}: 30, 20, then 40, 10.
\end{frame}
% ~~~~~~~~~~~~~~~~~~~~~~~~~~~~~~~~~~~~~~ A frm

% ~~~~~~~~~~~~~~~~~~~~~~~~~~~~~~~~~~~~~~ V
\subsection[Errors]{Errors}
% ~~~~~~~~~~~~~~~~~~~~~~~~~~~~~~~~~~~~~~ A




% ~~~~~~~~~~~~~~~~~~~~~~~~~~~~~~~~~~~~~~ V frm
\begin{frame}[t,fragile] \frametitle{Underflow: pop on an empty stack}
  \lstcoderightstep[0.5][basicstyle=\ttfamily\scriptsize][basicstyle=\ttfamily\tiny]{codeoutput/Underflow.java}{Java}{codeoutput/Underflow.txt}

  \vspace{0.2em}
  \small The exception names the problem (\texttt{underflow}) and the stack trace shows the line: \texttt{Underflow.java:6}.
\end{frame}
% ~~~~~~~~~~~~~~~~~~~~~~~~~~~~~~~~~~~~~~ A frm




% ~~~~~~~~~~~~~~~~~~~~~~~~~~~~~~~~~~~~~~ V frm
\begin{frame}[t,fragile] \frametitle{Overflow: push on a full stack}
  \lstcoderightstep[0.6][basicstyle=\ttfamily\scriptsize][basicstyle=\ttfamily\scriptsize]{codeoutput/Overflow.java}{Java}{codeoutput/Overflow.txt}

  \vspace{0.2em}
  \small Catching \texttt{IllegalStateException} lets the program continue. The stack is unchanged after the failed push.
\end{frame}
% ~~~~~~~~~~~~~~~~~~~~~~~~~~~~~~~~~~~~~~ A frm

% ~~~~~~~~~~~~~~~~~~~~~~~~~~~~~~~~~~~~~~ V
\hSection[Growing]{A Generic Stack That Grows}
% ~~~~~~~~~~~~~~~~~~~~~~~~~~~~~~~~~~~~~~ A




% ~~~~~~~~~~~~~~~~~~~~~~~~~~~~~~~~~~~~~~ V frm
\begin{frame}[t,fragile] \frametitle{\texttt{ArrayStack<T>}: any type, no fixed size}
  \begin{columns}[T]
    \begin{column}{0.49\textwidth}
      \lstcode[basicstyle=\ttfamily\tiny, lastline=20]{codeoutput/ArrayStack.java}{Java}
    \end{column}
    \begin{column}{0.49\textwidth}
      \lstcode[basicstyle=\ttfamily\tiny, firstline=21, firstnumber=21]{codeoutput/ArrayStack.java}{Java}
    \end{column}
  \end{columns}
  \small The \textbf{other} convention: the top is \texttt{a[size - 1]}. Full: \textbf{double}; one quarter full: \textbf{halve}.
  Java cannot make \texttt{new T[n]}: we make \texttt{Object[]} and cast.
\end{frame}
% ~~~~~~~~~~~~~~~~~~~~~~~~~~~~~~~~~~~~~~ A frm




% ~~~~~~~~~~~~~~~~~~~~~~~~~~~~~~~~~~~~~~ V frm
\begin{frame}[t,fragile] \frametitle{Watching the capacity}
  \lstcoderightstep[0.6][basicstyle=\ttfamily\scriptsize][basicstyle=\ttfamily\scriptsize]{codeoutput/Growing.java}{Java}{codeoutput/Growing.txt}

  \vspace{0.2em}
  \small The capacity doubles at sizes 1, 2, 4, 8 and halves when the stack is one quarter full: 16 $\to$ 8 at size 4, 8 $\to$ 4 at size 2, 4 $\to$ 2 at size 1.
\end{frame}
% ~~~~~~~~~~~~~~~~~~~~~~~~~~~~~~~~~~~~~~ A frm

% ~~~~~~~~~~~~~~~~~~~~~~~~~~~~~~~~~~~~~~ V
\hSection[Applications]{Applications}
% ~~~~~~~~~~~~~~~~~~~~~~~~~~~~~~~~~~~~~~ A

% ~~~~~~~~~~~~~~~~~~~~~~~~~~~~~~~~~~~~~~ V
\subsection[Reverse]{Reversing}
% ~~~~~~~~~~~~~~~~~~~~~~~~~~~~~~~~~~~~~~ A




% ~~~~~~~~~~~~~~~~~~~~~~~~~~~~~~~~~~~~~~ V frm
\begin{frame}[t,fragile] \frametitle{Reversing a word}
  \lstcoderightstep[0.64][basicstyle=\ttfamily\scriptsize][basicstyle=\ttfamily\scriptsize]{codeoutput/Reverse.java}{Java}{codeoutput/Reverse.txt}

  \vspace{0.2em}
  \small Push every letter, then pop them all: last in, first out reverses the order. Turkish letters are \texttt{char}s too.
\end{frame}
% ~~~~~~~~~~~~~~~~~~~~~~~~~~~~~~~~~~~~~~ A frm

% ~~~~~~~~~~~~~~~~~~~~~~~~~~~~~~~~~~~~~~ V
\subsection[Brackets]{Brackets}
% ~~~~~~~~~~~~~~~~~~~~~~~~~~~~~~~~~~~~~~ A




% ~~~~~~~~~~~~~~~~~~~~~~~~~~~~~~~~~~~~~~ V frm
\begin{frame}[t,fragile] \frametitle{Checking brackets}
  \lstcoderightstep[0.64][basicstyle=\ttfamily\tiny][basicstyle=\ttfamily\scriptsize]{codeoutput/Brackets.java}{Java}{codeoutput/Brackets.txt}

  \vspace{0.2em}
  \small \texttt{"([\{".indexOf(open) != ")]\}".indexOf(c)}: an opening and a closing bracket match when they are at the same position.
\end{frame}
% ~~~~~~~~~~~~~~~~~~~~~~~~~~~~~~~~~~~~~~ A frm

% ~~~~~~~~~~~~~~~~~~~~~~~~~~~~~~~~~~~~~~ V
\subsection[Postfix]{Postfix}
% ~~~~~~~~~~~~~~~~~~~~~~~~~~~~~~~~~~~~~~ A




% ~~~~~~~~~~~~~~~~~~~~~~~~~~~~~~~~~~~~~~ V frm
\begin{frame}[t,fragile] \frametitle{Evaluating postfix expressions}
  \lstcoderightstep[0.64][basicstyle=\ttfamily\tiny][basicstyle=\ttfamily\tiny]{codeoutput/Postfix.java}{Java}{codeoutput/Postfix.txt}

  \vspace{0.2em}
  \small The trace shows the stack after each token. \texttt{b} is popped \textbf{first}: for \texttt{3 -} that is the right operand.
\end{frame}
% ~~~~~~~~~~~~~~~~~~~~~~~~~~~~~~~~~~~~~~ A frm

% ~~~~~~~~~~~~~~~~~~~~~~~~~~~~~~~~~~~~~~ V
\subsection[Binary]{Number bases}
% ~~~~~~~~~~~~~~~~~~~~~~~~~~~~~~~~~~~~~~ A




% ~~~~~~~~~~~~~~~~~~~~~~~~~~~~~~~~~~~~~~ V frm
\begin{frame}[t,fragile] \frametitle{Decimal to binary}
  \lstcoderightstep[0.6][basicstyle=\ttfamily\tiny][basicstyle=\ttfamily\scriptsize]{codeoutput/Binary.java}{Java}{codeoutput/Binary.txt}

  \vspace{0.2em}
  \small Remainders are produced from the last digit to the first; the stack turns them around.
\end{frame}
% ~~~~~~~~~~~~~~~~~~~~~~~~~~~~~~~~~~~~~~ A frm

% ~~~~~~~~~~~~~~~~~~~~~~~~~~~~~~~~~~~~~~ V
\subsection[Undo]{Undo and redo}
% ~~~~~~~~~~~~~~~~~~~~~~~~~~~~~~~~~~~~~~ A




% ~~~~~~~~~~~~~~~~~~~~~~~~~~~~~~~~~~~~~~ V frm
\begin{frame}[t,fragile] \frametitle{Undo and redo with two stacks}
  \begin{columns}[T]
    \begin{column}{0.49\textwidth}
      \lstcode[basicstyle=\ttfamily\tiny, lastline=27]{codeoutput/Undo.java}{Java}
    \end{column}
    \begin{column}{0.49\textwidth}
      \lstcode[basicstyle=\ttfamily\tiny, firstline=28, firstnumber=28]{codeoutput/Undo.java}{Java}
      \uncover<2->{\lstoutput[basicstyle=\ttfamily\tiny]{codeoutput/Undo.txt}}
    \end{column}
  \end{columns}
\end{frame}
% ~~~~~~~~~~~~~~~~~~~~~~~~~~~~~~~~~~~~~~ A frm




% ~~~~~~~~~~~~~~~~~~~~~~~~~~~~~~~~~~~~~~ V frm
\begin{frame}[t,fragile] \frametitle{Undo and redo: what happened}
  \lstoutput[basicstyle=\ttfamily\footnotesize]{codeoutput/Undo.txt}
  \begin{itemize}\small
    \item Each stack holds whole texts: the state \emph{before} each change.
    \item \texttt{undo} moves the current text to \emph{redo}; \texttt{redo} moves it back to \emph{undo}.
    \item After \texttt{undo, undo, redo}, typing \texttt{" Java"} clears the redo stack: the \texttt{"!"} can never come back.
    \item Editors store \emph{changes} (commands) instead of whole texts, to save memory.
  \end{itemize}
\end{frame}
% ~~~~~~~~~~~~~~~~~~~~~~~~~~~~~~~~~~~~~~ A frm

% ~~~~~~~~~~~~~~~~~~~~~~~~~~~~~~~~~~~~~~ V
\subsection[Calls]{The call stack}
% ~~~~~~~~~~~~~~~~~~~~~~~~~~~~~~~~~~~~~~ A




% ~~~~~~~~~~~~~~~~~~~~~~~~~~~~~~~~~~~~~~ V frm
\begin{frame}[t,fragile] \frametitle{The call stack of the JVM}
  \lstcoderightstep[0.54][basicstyle=\ttfamily\tiny][basicstyle=\ttfamily\tiny]{codeoutput/CallStack.java}{Java}{codeoutput/CallStack.txt}

  \vspace{0.2em}
  \small Calls go in, returns come out in the reverse order. Without a base case the call stack overflows; the exact depth depends on the JVM and on \texttt{-Xss}.
\end{frame}
% ~~~~~~~~~~~~~~~~~~~~~~~~~~~~~~~~~~~~~~ A frm

% ~~~~~~~~~~~~~~~~~~~~~~~~~~~~~~~~~~~~~~ V
\hSection[Library]{Stacks from the Java Library}
% ~~~~~~~~~~~~~~~~~~~~~~~~~~~~~~~~~~~~~~ A




% ~~~~~~~~~~~~~~~~~~~~~~~~~~~~~~~~~~~~~~ V frm
\begin{frame}[t,fragile] \frametitle{\texttt{ArrayDeque} as a stack}
  \lstcoderightstep[0.56][basicstyle=\ttfamily\tiny][basicstyle=\ttfamily\tiny]{codeoutput/DequeDemo.java}{Java}{codeoutput/DequeDemo.txt}

  \vspace{0.2em}
  \small \texttt{push}, \texttt{pop}, \texttt{peek} work on the \textbf{front}. \texttt{poll} and \texttt{peek} return \texttt{null} on an empty deque; \texttt{pop} throws.
\end{frame}
% ~~~~~~~~~~~~~~~~~~~~~~~~~~~~~~~~~~~~~~ A frm




% ~~~~~~~~~~~~~~~~~~~~~~~~~~~~~~~~~~~~~~ V frm
\begin{frame}[t,fragile] \frametitle{The old \texttt{java.util.Stack}}
  \lstcoderightstep[0.6][basicstyle=\ttfamily\tiny][basicstyle=\ttfamily\tiny]{codeoutput/LegacyStack.java}{Java}{codeoutput/LegacyStack.txt}

  \vspace{0.2em}
  \small \hCode{Stack} prints bottom-first, \hCode{ArrayDeque} top-first. \texttt{search} counts from the top (1 = top). Being a \hCode{Vector}, \hCode{Stack} even lets you insert at the bottom: not a real stack.
\end{frame}
% ~~~~~~~~~~~~~~~~~~~~~~~~~~~~~~~~~~~~~~ A frm

% ~~~~~~~~~~~~~~~~~~~~~~~~~~~~~~~~~~~~~~ V
\hSection[Summary]{Summary}
% ~~~~~~~~~~~~~~~~~~~~~~~~~~~~~~~~~~~~~~ A




% ~~~~~~~~~~~~~~~~~~~~~~~~~~~~~~~~~~~~~~ V frm
\begin{frame}[t,fragile] \frametitle{What we saw}
  \begin{center}\small
  \begin{tabular}{lL{0.28\textwidth}L{0.36\textwidth}}
    \toprule
    \textbf{Program}       & \textbf{Shows}                       & \textbf{Key line} \\
    \midrule
    \texttt{IntStack}      & fixed array, \texttt{top}            & \texttt{a[++top] = x}, \texttt{return a[top--]} \\
    \texttt{ArrayStack}    & generic, grows and shrinks           & \texttt{if (size == a.length) resize(2 * a.length)} \\
    \texttt{Reverse}       & LIFO reverses                        & push all, pop all \\
    \texttt{Brackets}      & matching pairs                       & pop must match the closing bracket \\
    \texttt{Postfix}       & evaluating expressions               & pop \texttt{b}, pop \texttt{a}, push \texttt{a op b} \\
    \texttt{Binary}        & reversing digits                     & push \texttt{n \% 2} \\
    \texttt{Undo}          & two stacks                           & a new edit clears \emph{redo} \\
    \texttt{CallStack}     & the JVM's own stack                  & \texttt{StackOverflowError} \\
    \texttt{ArrayDeque}    & the library stack                    & \texttt{Deque<T> s = new ArrayDeque<>()} \\
    \bottomrule
  \end{tabular}
  \end{center}
\end{frame}
% ~~~~~~~~~~~~~~~~~~~~~~~~~~~~~~~~~~~~~~ A frm




% ~~~~~~~~~~~~~~~~~~~~~~~~~~~~~~~~~~~~~~ V frm
\begin{frame}[t,fragile] \frametitle{Exercises}
  \begin{enumerate}\small
    \item Add a method \texttt{clear()} to \texttt{IntStack}. Why is \texttt{top = -1} enough?
    \item Change \texttt{IntStack} to the \texttt{size} convention. Which lines change?
    \item In \texttt{Growing}, push 100 elements. How many times is the array copied? How many elements are copied in total?
    \item Extend \texttt{Postfix} with \texttt{\%} and \texttt{\textasciicircum} (power). What happens with \texttt{"1 +"}?
    \item Make \texttt{Brackets} ignore brackets inside string literals \texttt{"..."}.
    \item Rewrite \texttt{Undo} with \texttt{ArrayStack<String>}. Which methods do you need?
  \end{enumerate}
\end{frame}
% ~~~~~~~~~~~~~~~~~~~~~~~~~~~~~~~~~~~~~~ A frm




% ~~~~~~~~~~~~~~~~~~~~~~~~~~~~~~~~~~~~~~ V frm
\begin{frame}[t,fragile]
	\begin{center}
		\vspace{4cm}
		\hRemark{\Huge Questions?}
	\end{center}
\end{frame}
% ~~~~~~~~~~~~~~~~~~~~~~~~~~~~~~~~~~~~~~ A frm



\end{document}
